\section*{Chapitre \ref{ch:b3:petrological-thermodynamics} --- Thermodynamique des roches}

\begin{solution}{exo:b3:petrological-thermodynamics:1}
$\Delta_rH = -3029.3 - 910.7 + 3935.0 = \qty{-5.0}{kJ}$~; $\Delta_rS = 133.5 +
41.5 - 207.4 = \qty{-32.4}{J.K^{-1}}$~; $\Delta_rV = 60.40 + 22.69 - 100.07 =
\qty{-16.98}{cm^3} = \qty{-1.698}{J.bar^{-1}}$. Pente $32.4/1.698 =
\qty{19.1}{bar/K}$. À \qty{873}{K}~: $P = 1 + (873.15\times(-32.4) +
5000)/(-1.698) = \qty{13.7}{kbar}$.
\end{solution}

\begin{solution}{exo:b3:petrological-thermodynamics:2}
$c = 2$ et $p = 3 > c$~: la règle est violée. Les trois phases ne coexistent
que sur la droite univariante de \ce{Mg2SiO4 + SiO2 <=> 2 MgSiO3}~; une roche
prélevée au hasard ne s'y trouve pas, donc l'assemblage témoigne d'une réaction
incomplète, par exemple une forstérite entourée d'une couronne d'enstatite qui
l'a isolée du quartz.
\end{solution}

\begin{solution}{exo:b3:petrological-thermodynamics:3}
$a_{\text{grossulaire}} = 0.20^3 = 0.008$~; $a_{\text{pyrope}} = 0.15^3 =
\num{3.4e-3}$~; $a_{\text{anorthite}} = 0.40$.
\end{solution}

\begin{solution}{exo:b3:petrological-thermodynamics:4}
$T_c = 15\,000/(2\times8.314) = \qty{902}{K}$~; $W = 2\times8.314\times1000
= \qty{16.6}{kJ.mol^{-1}}$.
\end{solution}

\begin{solution}{exo:b3:petrological-thermodynamics:5}
$\Delta_rG = -5432.2 + 2\times(-2088.8) - 730.2 + 3\times3478.1 =
\qty{94.3}{kJ}$~; $P = 1 + 94\,300/6.61 = \qty{14.3}{kbar}$. La droite
expérimentale donne $22.80\times726.85 - 1093 = \qty{15.5}{kbar}$, soit
\qty{1.2}{kbar} de plus.
\end{solution}

\begin{solution}{exo:b3:petrological-thermodynamics:6}
$RT = \qty{7483}{J.mol^{-1}}$. En $x = 0.5$~: $7483\ln0.5 + 20\,000/4 = -5187 +
5000 = \qty{-187}{J.mol^{-1}}$. En $x = 0.111$~: $7483\times(0.111\ln0.111 +
0.889\ln0.889) + 20\,000\times0.111\times0.889 = -2609 + 1974 =
\qty{-635}{J.mol^{-1}}$, et la même valeur en 0.889 par symétrie. Des quantités
égales des deux flancs ont pour composition globale 0.5 et pour énergie de Gibbs
\qty{-635}{J.mol^{-1}}~: la séparation l'abaisse de \qty{448}{J.mol^{-1}}.
\end{solution}

\begin{solution}{exo:b3:petrological-thermodynamics:7}
$c + 2 = 5$ phases et 5 droites. (D) $A + C = 2B$~: $0 + 2/3 = 2\times1/3$. (A)
$B + D = 2C$~: $1/3 + 1 = 2\times2/3$. (C) $2A + D = 3B$~: $0 + 1 = 3\times1/3$.
(B) $A + 2D = 3C$~: $0 + 2 = 3\times2/3$.
\end{solution}

\begin{solution}{exo:b3:petrological-thermodynamics:8}
D'après la règle des segments inverses, la phase en 0.889 représente $(0.3 - 0.111)/(0.889 - 0.111) =
0.243$, soit \qty{24}{\percent}, et la phase en 0.111 \qty{76}{\percent}.
\end{solution}

\begin{solution}{exo:b3:petrological-thermodynamics:9}
$\partial P/\partial X_{\text{anorthite}} = -3RT/(6.21X_{\text{anorthite}}) =
-3\times8.314\times900/(6.21\times0.30) = \qty{-12000}{bar}$ par unité de
$X_{\text{anorthite}}$~: une erreur de 0.01 déplace la pression de \qty{120}{bar}.
\end{solution}

\begin{solution}{exo:b3:petrological-thermodynamics:10}
$G'_{\text{mix}} = RT\ln[x/(1 - x)] + W(1 - 2x)$ et $G''_{\text{mix}} =
RT/[x(1 - x)] - 2W$, nulle lorsque $x(1 - x) = RT/(2W)$. Pour $x = 0.3$~: $T =
2\times20\,000\times0.21/8.314 = \qty{1010}{K}$~; le \omterm{def:b3:petrological-thermodynamics:solvus}{solvus} est atteint à
$20\,000\times0.4/(8.314\ln(7/3)) = \qty{1136}{K}$. Entre les deux, le
cristal est métastable~: la séparation abaisserait son énergie, mais une
petite fluctuation l'élève, de sorte que la seconde phase doit germer.
\end{solution}

\begin{solution}{exo:b3:petrological-thermodynamics:11}
Puisque $G_{\text{mix}}(x) = G_{\text{mix}}(1 - x)$, la tangente commune touche
la courbe en deux points symétriques $x$ et $1 - x$ et est horizontale~:
$G'_{\text{mix}}(x) = 0$, qui est l'équation. En la résolvant, $T(x) = W(1 -
2x)/[R\ln((1 - x)/x)]$. Avec $x = 1/2 - \varepsilon$, $1 - 2x = 2\varepsilon$
et $\ln((1/2 + \varepsilon)/(1/2 - \varepsilon)) \approx 4\varepsilon$, donc
$T \to W/(2R)$ lorsque $\varepsilon \to 0$~; et pour $x < 1/2$, $T(x)$ croît
à mesure que $x$ se rapproche de 1/2, si bien que cette limite est le maximum.
\end{solution}

\begin{solution}{exo:b3:petrological-thermodynamics:12}
$(2.0 - 0.08\times10)/0.92 = \qty{1.30}{wt\percent}$. Le manganèse favorise
le grenat~: avec moins de manganèse dans la \omterm{def:b3:petrological-thermodynamics:pseudosection}{composition globale effective}, les
bordures croissent dans un autre champ et le long d'\omterm{def:b3:petrological-thermodynamics:pseudosection}{isoplèthes} différentes de
celles que prévoit l'analyse de la roche totale.
\end{solution}

\begin{solution}{pb:b3:petrological-thermodynamics:1}
\textbf{1.} $\Delta_rH^\circ = -6640.0 + 2\times(-2593.8) - 910.7 + 3\times4234.0
= \qty{-36.3}{kJ}$~; $\Delta_rS^\circ = 260.1 + 165.6 + 41.5 - 597.9 =
\qty{-130.7}{J.K^{-1}}$~; $\Delta_rV = 125.28 + 88.30 + 22.69 - 302.37 =
\qty{-66.10}{cm^3} = \qty{-6.61}{J.bar^{-1}}$.

\textbf{2.} $-36\,300 + 298.15\times130.7 = \qty{+2.67}{kJ}$~: positive, donc
l'anorthite est stable à la surface.

\textbf{3.} $P = 1 + (900\times(-130.7) + 36\,300)/(-6.61) = \qty{12.3}{kbar}$~;
pente $130.7/6.61 = \qty{19.8}{bar/K}$.

\textbf{4.} $1 + 81\,500/6.61 = \qty{12.3}{kbar}$ à \qty{900}{K}~; $1 +
119\,800/6.61 = \qty{18.1}{kbar}$ à \qty{1200}{K}.

\textbf{5.} À 626.85 et \qty{926.85}{\celsius}, la droite expérimentale donne
13.2 et \qty{20.0}{kbar}~: soit 0.9 et \qty{1.9}{kbar} au-dessus des tables.

\textbf{6.} $1914\times6.61 = \qty{12.7}{kJ}$. L'incertitude combinée est
$\sqrt{3.2^2 + (3\times2.0)^2 + (2\times2.0)^2 + 1.0^2} = \qty{8.0}{kJ}$~: l'écart
vaut 1.6 fois celle-ci, ce qui est possible pour une calorimétrie, et c'est la
raison pour laquelle la droite est tirée des expériences.

\textbf{7.} $\Delta_rC_p$ est faible parmi les solides, et ses deux intégrales
se compensent en partie~: les tables déplacent la droite de seulement 26 et \qty{112}{bar}.

\textbf{8.} \ce{SiO2}, \ce{Al2O3}, \ce{FeO}, \ce{MgO}, \ce{MnO}, \ce{CaO},
\ce{Na2O}, \ce{K2O} et \ce{H2O}~: $c = 9$. Phases~: grenat, biotite,
muscovite, plagioclase, disthène, quartz et fluide, $p = 7$.

\textbf{9.} $F = 9 - 7 + 2 = 4$~: pression, température et deux variables de composition
variables. $p = 7 \le c = 9$~: la règle est respectée, et l'assemblage peut être stable
sur un domaine du plan pression--température.

\textbf{10.} Deux polymorphes purs \ce{Al2SiO5} en équilibre se situent sur la
droite univariante disthène = sillimanite~: une relation entre $P$ et $T$, et un
géothermomètre donne alors la pression.

\textbf{11.} Le fluide impose le potentiel chimique de l'eau~: 8 composants libres
pour 6 phases solides, $F = 8 - 6 + 2 = 4$, et $p \le c$ reste
vérifiée.

\textbf{12.} $a_{\text{grossulaire}} = 0.10^3 = 0.001$, $a_{\text{anorthite}} =
0.30$; $\ln K = \ln(0.001/0.30^3) = -3.30$.

\textbf{13.} $K_D = (0.12/0.74)/(0.45/0.55) = 0.1622/0.8182 = 0.198$.

\textbf{14.} $T = (52\,635 + 2488)/(20.20 + 24.94\times1.618) = 55\,123/60.57
= \qty{910}{K}$, soit \qty{637}{\celsius}.

\textbf{15.} $P_0 = 22.80\times636.9 - 1093 = \qty{13429}{bar}$; $(3\times8.314
\times910.1/6.21)\ln(1/3) = \qty{-4016}{bar}$; $P = \qty{9.41}{kbar}$.

\textbf{16.} À \qty{9.41}{kbar}, \qty{644.2}{\celsius}, puis
\qty{9.55}{kbar}~; à cette pression, \qty{644.9}{\celsius} et
\qty{9.56}{kbar}, sans évolution ensuite~: \qty{645}{\celsius} et
\qty{9.56}{kbar}.

\textbf{17.} \qty{641}{\celsius} et \qty{8.79}{kbar}. Les tables placent la
courbe des pôles purs environ \qty{1}{kbar} trop bas (partie I)~; la droite
expérimentale est conservée.

\textbf{18.} Avec $X_{\text{an}} = 0.35$~: \qty{642}{\celsius} et
\qty{8.93}{kbar}, soit \qty{0.63}{kbar} de moins.

\textbf{19.} À \qty{918}{K}~: $P = 1 + (918\times12.6 - 7700)/0.571 =
\qty{6.8}{kbar}$, inférieure à \qty{9.56}{kbar}~: disthène, comme observé.

\textbf{20.} $z = \num{9.56e8}/(2800\times9.81) = \qty{34.8}{km}$.

\textbf{21.} $\Delta P = (3\times8.314\times918/6.21)\ln(0.07/0.10) =
\qty{-1.3}{kbar}$~: le cœur s'est formé environ \qty{4.8}{km} plus près de la surface. Le grenat
a crû pendant l'enfouissement, et sa bordure enregistre le point le plus profond.

\textbf{22.} Le schiste de Varden a atteint une pression maximale de
\textbf{\qty[text-series-to-math=true]{9.6}{kbar}} à
\qty{645}{\celsius}, pour un enfouissement d'environ \qty{35}{km}.
\end{solution}
